\SourceLeaf{211}%
\noindent gesetzt werden kann. Hier wird also
\[
\omega=aa\sin u^2\,\mathrm{d}t^2+(aa\cos u^2+bb\sin u^2)\mathrm{d}u^2
\]
\noindent und die Differentialformel \(\omega=0\) gibt, wenn wir Kürze halber \(\BareRoot{(1-\frac{bb}{aa})}=\varepsilon\) setzen (insofern die Revolutionshalbaxe \(b<a\)),
\[
0=\mathrm{d}t\mp i\,\mathrm{d}u\mathord{.}\BareRoot{(\mathrm{cotang}\,u^2+1-\varepsilon\varepsilon)}
\]

Setzt man hier
\[
\BareRoot{(1-\varepsilon\varepsilon)}\mathord{.}\mathrm{tang}\,u=\mathrm{tang}\,w
\]
\noindent wo, bei der Anwendung auf das Erdsphäroid, \(90^0-w\) die geographische Breite und \(t\) die Länge vorstellen wird, so verwandelt sich diese Gleichung in
\[
0=\mathrm{d}t\mp i\,\mathrm{d}w\mathord{.}\frac{1-\varepsilon\varepsilon}{(1-\varepsilon\varepsilon\cos w^2)\sin w}
\]
\noindent deren Integration
\[
\text{Const.}=t\pm i\log\mathord{.}\Bigl\{\mathrm{cotang}\tfrac{1}{2}w\mathord{.}\Bigl(\frac{1-\varepsilon\cos w}{1+\varepsilon\cos w}\Bigr)^{\frac{1}{2}\varepsilon}\Bigr\}
\]
\noindent gibt. Man hat daher, indem \(f\) eine willkürliche Function bedeutet, für \(X\) den reellen und für \(iY\) den imaginären Theil von
\[
f\Bigl(t+i\log\Bigl\{\mathrm{cotang}\tfrac{1}{2}w\mathord{.}\Bigl(\frac{1-\varepsilon\cos w}{1+\varepsilon\cos w}\Bigr)^{\frac{1}{2}\varepsilon}\Bigr\}\Bigr)
\]
\noindent zu nehmen. --- Wählt man für \(f\) eine lineäre Function, d. i. \(f\upsilon=k\upsilon\), so wird
\[
X=kt,\quad Y=k\log\mathrm{cotang}\tfrac{1}{2}w-\tfrac{1}{2}k\varepsilon\log\frac{1+\varepsilon\cos w}{1-\varepsilon\cos w}
\]
\noindent welches eine der \textsc{Mercator}schen analoge Projection gibt.

Nimmt man hingegen für \(f\) eine imaginäre Exponentialfunction \(f\upsilon=ke^{i\lambda\upsilon}\), so wird
\[
\begin{aligned}
X&=k\mathord{.}\mathrm{tang}\tfrac{1}{2}w^{\lambda}\mathord{.}\Bigl(\frac{1+\varepsilon\cos w}{1-\varepsilon\cos w}\Bigr)^{\frac{1}{2}\varepsilon\lambda}\mathord{.}\cos\lambda t,\\
Y&=k\,\mathrm{tang}\tfrac{1}{2}w^{\lambda}\mathord{.}\Bigl(\frac{1+\varepsilon\cos w}{1-\varepsilon\cos w}\Bigr)^{\frac{1}{2}\varepsilon\lambda}\mathord{.}\sin\lambda t
\end{aligned}
\]
\noindent welches, wenn man \(\lambda=1\) setzt, eine der stereographischen Polarprojection analoge, und allgemein, eine zur Darstellung eines Theils der Erdoberfläche, insofern man auf die Abplattung Rücksicht nehmen soll, sehr zweckmässige Projection gibt.

Was über den andern Fall, wo \(b>a\) ist, zu sagen ist, lässt sich zwar leicht aus dem vorhergehenden unmittelbar ableiten, wo, wenn man dieselben
